\chapter{Memory and State Persistence}
\section{Four forms of continuity}
A chat transcript and an execution checkpoint solve different problems. The transcript records what participants said. A checkpoint records what the runtime has actually committed, including pending or completed actions. RAG provides separately managed knowledge; it is not a hidden copy of conversation memory.
\begin{center}\begin{tikzpicture}[scale=.95,transform shape]
\node[darkbox,text width=4cm] (root) at (0,5.4) {META-HARNESS};
\node[box] (conv) at (-4,3.7) {Conversation history\\messages + attachments};
\node[box] (task) at (0,3.7) {Task state\\checkpoints + operations};
\node[box] (rag) at (4,3.7) {External knowledge\\devices + document chunks};
\node[box] (working) at (0,1.9) {Private working context\\delta / memo / compaction};
\node[darkbox,text width=6.8cm] (pg) at (-2,.2) {Core PostgreSQL\\transcript + private windows + task ledgers};
\node[darkbox] (rpg) at (4,.2) {External RAG\\pgvector + original files};
\draw[arr] (root)--(conv);\draw[arr] (root)--(task);\draw[arr] (root)--(rag);
\draw[arr] (conv)--(pg);\draw[arr] (task)--(working);\draw[arr] (working)--(pg);\draw[arr] (rag)--(rpg);
\draw[arr,dashed] (rag.west) -- node[lab,above]{MCP evidence} (working.east);
\end{tikzpicture}
\captionof{figure}{Memory ownership. Chat and task state share core PostgreSQL. External knowledge remains in the independent RAG database and file volumes.}\end{center}
\begin{description}
\item[Conversation memory] Ordered user/assistant messages, automatically derived chat titles, attachment bytes and compact tool traces. Saved messages remain visible even when model context is summarized.
\item[Working memory] The private input window supplied to the selected model: relevant messages, decisions, observations and opaque provider compaction/reasoning items. This is separate from UI text.
\item[Persistent task state] Goal, immutable plan/criteria, current and completed steps, pending action, operation ledger, approvals, retries, verification evidence, context cursors, budget baselines and lease information.
\item[Long-term domain knowledge] Device/issue data and searchable document passages in RAG. The agent retrieves relevant evidence through permitted MCP calls; it does not learn every uploaded manual permanently.
\end{description}
\section{Conversation reconstruction}
\texttt{ChatsService} inserts messages and attachments transactionally, serializes response generation per chat and links a reply to its originating user message. The unique reply linkage prevents multiple final replies to the same turn. Ordinary model context has a model identifier, complete private input items and a message-position watermark. Only new input is added when that saved window is compatible; model changes rebuild from visible input.

Attachments are stored as PostgreSQL \texttt{bytea} for downloading. Ordinary attachment content is not extracted into LLM context. Uploading a PDF with the chat plus button is therefore different from uploading it into RAG Documents.
\source{\path{core/src/chats/chats.service.ts}; \path{core/src/llm/context-window.ts}: \texttt{ChatContextWindow}; \path{database/schema.sql}; \path{ui/src/components/chat/attachment-card.tsx}.}
\nextpage{What PostgreSQL actually saves}
\section*{The core persistence model}
Core uses parameterized \texttt{pg} SQL, not an ORM. The schema snapshot and ordered migrations are applied inside a transaction with a database advisory lock. \texttt{schema\_migrations} stores a checksum for each upgrade; changing an applied migration is a startup error rather than an implicit schema rewrite.
\begin{tabularx}{\linewidth}{@{}lX@{}}\toprule
Entity & Persisted purpose\\\midrule
\texttt{chats} & Title, activity timestamps and private conversation context\\
\texttt{messages} & Role/text, order, reply linkage, tool traces and continuation state\\
\texttt{attachments} & Filename, MIME metadata, size and original bytes\\
\texttt{mcp\_servers} & Endpoints, discovered tools, permissions, status and credential-variable reference\\
\texttt{agent\_tasks} & Request/goal/plan/criteria, progress, next action, private context and budgets\\
\texttt{task\_steps} & Ordered decisions and execution step status\\
\texttt{task\_operations} & Stable operation identity, action, policy digest, dispatch/result/error and attempts\\
\texttt{task\_approvals} & Exact operation/action, policy identity and operator decision\\
\texttt{task\_verifications} & Criterion ID, pass/fail and actual observed evidence\\
\texttt{task\_events} & Sequenced progress/model/error/checkpoint events\\
\texttt{task\_schedules} & Cron/one-time schedule, timezone and task template\\\bottomrule
\end{tabularx}
Historical \texttt{task\_triggers} and \texttt{task\_trigger\_events} tables remain in the immutable schema but are inactive. Secrets are absent from these tables: profiles point to names in root \texttt{.env}.
\section*{State is saved before it becomes external work}
The worker first commits a model decision as \texttt{next\_action}. It then persists an operation identity and approval requirement. Immediately before sending a permitted external request, it commits \texttt{DISPATCHING}. Actual results, verification and progress are separate transactions. This ordering allows recovery to distinguish an unsent proposal from an action whose outcome is uncertain.

\texttt{fenced()} locks the task row and verifies RUNNING status plus the current lease-owner UUID. A stale worker cannot overwrite a pause/cancel or another owner's checkpoint. A session advisory lock supplies additional exclusion even if lease time expires while a process is still alive.
\section*{Schema evolution}
Migration 001 adds durable tasks, 002 adds estimated-cost accounting and 003 adds context-delivery state, prompt-token/turn/active-time baselines and a continuation counter. Task-operation evidence stays in PostgreSQL across all three. Foreign keys and cascade rules prevent orphaned task history when its conversation is deleted.
\source{\path{core/src/database/database.service.ts}: \texttt{onModuleInit()}, \texttt{transaction()}; \path{database/migrations/001_durable_tasks.sql}; \path{database/migrations/002_task_cost_budget.sql}; \path{database/migrations/003_task_context_checkpoints.sql}; \path{core/src/tasks/task-worker.ts}: \texttt{fenced()}.}
\nextpage{Incremental context and long-task continuity}
\section*{Why full snapshots were replaced}
The old task planner appended recent chat plus every full tool result on each turn. The private window grew with repeated evidence, while cumulative token limits could stop a task before its final source check. The current design separates authoritative storage from what the next model prompt needs.

\texttt{taskContextDelta()} fingerprints each operation's status, result, error and attempt count. A new prompt receives recent conversation and current tool schemas once. Later turns receive only changed operation observations and changed schemas, plus a bounded ledger and deterministic criterion statuses. Fingerprints and the resulting private window are committed together: a crash cannot mark evidence as delivered without saving the context that received it.
\section*{Bounded results with exact evidence access}
Each changed result is sent in full up to 12,000 characters; larger results use a marked preview. An update targets at most 64,000 characters of changed evidence, with pending changes delivered later. The ledger includes the latest forty operations; completed-step labels include the latest ten. Full original results remain in storage.

If exact detail is needed, the model chooses an internal \texttt{evidence} action. \texttt{readTaskEvidence()} resolves only operation IDs belonging to that task and applies a JSON pointer plus bounded offset/limit. It returns at most fifty array items or 12,000 text characters, also bounding large array pages. A nested oversized item can be inspected through a more specific pointer. This is neither an external replay nor a new proof of external execution.
\section*{Two distinct summarization boundaries}
\begin{enumerate}
\item Before inference, core counts the current model input. At 90\% of configured context capacity, the selected model's Responses compaction endpoint produces a canonical replacement window. If compaction remains too large or unavailable, original evidence is retained and an error is surfaced.
\item At a task prompt's turn/token/active-time budget, \texttt{summarizeTaskContext()} asks the same selected model for a factual working memo, with a 2,048-token output bound. \texttt{continuePrompt()} commits that memo, refreshes budget baselines and automatically requeues the task.
\end{enumerate}
The memo preserves identifiers, goals, sources, constraints, uncertain/declined writes and next steps. Immutable criteria and operation identities remain in database fields rather than depending solely on summary prose. Legacy windows receive one checkpoint when resumed. A failed summary does not replace the old window.
\note{What remains limited}{Summaries can omit facts or lose nuance. Exact stored-evidence reads mitigate loss but do not guarantee perfect model recall. Tokens/cost are cumulative; cost and continuation caps pause execution. Interrupted provider requests may be billed even when usage is unavailable. No semantic search over chat memory or learned personal-memory store is implemented.}
\source{\path{core/src/tasks/task-context.ts}: \texttt{taskContextDelta()}, \texttt{readTaskEvidence()}; \path{core/src/llm/llm.service.ts}: \texttt{summarizeTaskContext()}; \path{core/src/llm/context-window.ts}: \texttt{countContext()}, \texttt{compactContext()}.}
